\chapter{A integral de Lebesgue}\label{ch:b3:lebesgue}

A integral de Riemann fatia o \emph{domínio} em pequenos intervalos; a
de Lebesgue fatia a \emph{imagem}: para integrar $f$, meça os
conjuntos $\{f > t\}$. A mudança parece inocente e é revolucionária. Limites e
integrais, eternamente em desavença na teoria de Riemann (exigia-se
convergência uniforme!), reconciliam-se por três teoremas de
convergência --- convergência monótona, Fatou, convergência dominada
--- cujas hipóteses são de uma fraqueza quase embaraçosa. Este capítulo
constrói a integral sobre um \omterm{def:b3:measure:measure}{espaço de medida} arbitrário $(X, \mathcal A,
\mu)$,
demonstra os três teoremas, esclarece a relação exata com a integral de
Riemann (uma função limitada é \omterm{def:b3:lebesgue:l1}{Riemann-integrável} se, e somente se, é
\omterm{def:b3:topology:continuity}{contínua} \omterm{def:b3:lebesgue:l1}{quase em toda parte}) e industrializa a derivação de integrais
dependentes de parâmetro --- a técnica que o problema de fim de semana
usa para calcular $\int_0^\infty\frac{\sin x}x\,\dd x$ e $\int_\R \eu^{-x^2}\dd x$.

\section{Funções mensuráveis}

\begin{definition}\label{def:b3:lebesgue:measurable}
Sejam $(X, \mathcal A)$ e $(Y, \mathcal B)$ espaços mensuráveis. $f \colon X \to Y$ é
\emph{mensurável}\index{função mensurável} se $f^{-1}(B) \in \mathcal A$ para todo $B \in
\mathcal B$. Para
funções reais (ou a valores em $[-\infty,+\infty]$), $Y = \R$ carrega sua
$\sigma$-álgebra de Borel, e basta verificar $f^{-1}(\intoo t{+\infty}) = \{f > t\} \in
\mathcal A$ para todo $t \in \R$: os
bons conjuntos $\{B :
f^{-1}(B) \in \mathcal A\}$ formam uma $\sigma$-álgebra (as pré-imagens
comutam com as operações de conjuntos) que contém as semirretas
geradoras (\cref{def:b3:measure:borel}, \cref{met:b3:measure:goodsets}).
\end{definition}

\begin{proposition}\label{prop:b3:lebesgue:stability}
(a) Composições de aplicações \omterm{def:b3:lebesgue:measurable}{mensuráveis} são \omterm{def:b3:lebesgue:measurable}{mensuráveis}; \omterm{def:b3:topology:continuity}{aplicações
contínuas} são \omterm{def:b3:lebesgue:measurable}{Borel-mensuráveis}.
(b) Se $f, g \colon X \to \R$ são \omterm{def:b3:lebesgue:measurable}{mensuráveis}, também o são $f + g$, $fg$, $\max(f,g)$,
$\abs f$, $\lambda f$.
(c) Se as $(f_n)$ são \omterm{def:b3:lebesgue:measurable}{mensuráveis} com valores em $[-\infty, +\infty]$, então $\sup_nf_n$,
$\inf_nf_n$, $\limsup f_n$, $\liminf f_n$ são \omterm{def:b3:lebesgue:measurable}{mensuráveis}; se $f_n \to f$ pontualmente, $f$ é
\omterm{def:b3:lebesgue:measurable}{mensurável}.
\end{proposition}

\begin{proof}
(a) $(g\circ f)^{-1}(B) = f^{-1}(g^{-1}(B))$; a \omterm{def:b3:topology:continuity}{continuidade} dá a \omterm{def:b3:lebesgue:measurable}{mensurabilidade} via os abertos geradores
(\cref{pb:b3:measure:1}, questão 10, em forma geral).
(b) $(f, g) \colon X \to \R^2$ é \omterm{def:b3:lebesgue:measurable}{mensurável} para a $\sigma$-álgebra de Borel de $\R^2$
--- verifique nos blocos abertos, que a geram (os abertos de $\R^2$
são reuniões enumeráveis de blocos racionais): $(f,g)^{-1}(U\times V) = f^{-1}(U)\cap g^{-1}(V)$ --- e $+, \times, \max$ são
\omterm{def:b3:topology:continuity}{contínuas} $\R^2 \to \R$: componha.
(c) $\{\sup f_n > t\} = \bigcup_n\{f_n > t\}$; $\inf = -\sup(-)$; $\limsup = \inf_N\sup_{n \geq N}$; um limite pontual é seu próprio $\limsup$.
\end{proof}

\begin{definition}\label{def:b3:lebesgue:simple}
Uma \emph{função simples}\index{função simples} é uma \omterm{def:b3:lebesgue:measurable}{função mensurável}
com um número finito de valores: $s = \sum_{i=1}^n
c_i\,\mathbf 1_{A_i}$, com $A_i \in \mathcal A$ disjuntos e $c_i \geq
0$ (para
a teoria não negativa). Sua integral é
\[
\int s\,\dd\mu = \sum_i c_i\,\mu(A_i) \in [0, +\infty]
\]
(com a convenção $0\cdot\infty = 0$); o valor não depende da \omterm{def:b3:representations:rep}{representação} (refine
duas partições).
\end{definition}

\begin{theorem}[Aproximação por funções simples]
\label{thm:b3:lebesgue:approximation}
Toda \omterm{def:b3:lebesgue:measurable}{função mensurável} $f \colon X \to [0, +\infty]$ é o limite pontual de uma sequência
\emph{crescente} de \omterm{def:b3:lebesgue:simple}{funções simples}:
\[
s_n = \sum_{k=1}^{n2^n} \frac{k-1}{2^n}\,
\mathbf 1_{\{\frac{k-1}{2^n} \leq f < \frac k{2^n}\}}
+ n\,\mathbf 1_{\{f \geq n\}} \nearrow f .
\]
\end{theorem}

\begin{proof}
Cada $s_n$ é simples (os conjuntos são pré-imagens de borelianos).
Monotonicidade: passar de $n$ a $n+1$ divide cada nível diádico em
dois e nunca diminui o valor atribuído (um ponto com $\frac{k-1}{2^n} \leq f(x) < \frac k{2^n}$ recebe $\frac{2k-2}{2^{n+1}}$
ou $\frac{2k-1}{2^{n+1}}$, ambos $\geq
\frac{k-1}{2^n}$; e o teto $n$ também sobe). Convergência: se
$f(x) < \infty$, para $n > f(x)$ temos $f(x) - s_n(x) \leq
2^{-n}$; se $f(x) = \infty$, $s_n(x) = n \to \infty$.
\end{proof}

\section{A integral e os teoremas de convergência}

\begin{definition}\label{def:b3:lebesgue:integral}
Para $f \geq 0$ \omterm{def:b3:lebesgue:measurable}{mensurável}:
\[
\int f \,\dd\mu = \sup\Bigl\{\int s\,\dd\mu : s \text{ simples},
\ 0 \leq s \leq f\Bigr\} \in [0, +\infty].
\]
Ela é monótona em $f$ por construção e estende o caso simples (para
$f$ simples, o supremo é atingido em $f$: comparação de integrais
de \omterm{def:b3:lebesgue:simple}{funções simples} via refinamentos comuns).
\end{definition}

\begin{theorem}[Convergência monótona, Beppo Levi]
\label{thm:b3:lebesgue:mct}
Se $0 \leq f_n \nearrow f$ pontualmente (com todas \omterm{def:b3:lebesgue:measurable}{mensuráveis}), então
\[
\int f_n\,\dd\mu \nearrow \int f\,\dd\mu .
\]
\index{teorema da convergência monótona}
\end{theorem}

\begin{proof}
$f$ é \omterm{def:b3:lebesgue:measurable}{mensurável} (\cref{prop:b3:lebesgue:stability}(c)) e $\int f_n$ cresce até algum $L \leq \int f$
(monotonicidade). Reciprocamente, fixe uma \omterm{def:b3:lebesgue:simple}{função simples} $s = \sum c_i\mathbf 1_{A_i} \leq f$ e $\theta \in (0,1)$;
os conjuntos $E_n = \{f_n \geq \theta s\}$ são \omterm{def:b3:lebesgue:measurable}{mensuráveis} e crescem até $X$ (onde $s(x) > 0$:
$f(x)
\geq s(x) > \theta s(x)$, de modo que, a partir de certo índice, $f_n(x) \geq \theta
s(x)$; onde $s(x) = 0$:
trivialmente). Então
\[
\int f_n \geq \int_{E_n}\theta s\,\dd\mu
= \theta\sum_i c_i\,\mu(A_i \cap E_n)
\xrightarrow[n\to\infty]{} \theta\sum_ic_i\,\mu(A_i)
= \theta\int s
\]
por \omterm{def:b3:topology:continuity}{continuidade} por baixo (\cref{prop:b3:measure:basics}(c)). Logo $L \geq \theta\int s$ para todo
$\theta < 1$ e toda $s
\leq f$ simples: $L \geq \int f$.
\end{proof}

\begin{corollary}\label{cor:b3:lebesgue:additivity}
Para $f, g \geq 0$ \omterm{def:b3:lebesgue:measurable}{mensuráveis} e $c \geq 0$: $\int(f + g) =
\int f + \int g$ e $\int cf = c\int f$; e, para uma série de
\omterm{def:b3:lebesgue:measurable}{funções mensuráveis} não negativas, $\int\sum_nf_n =
\sum_n\int f_n$.
\end{corollary}

\begin{proof}
Para \omterm{def:b3:lebesgue:simple}{funções simples}, a aditividade é um cálculo em um refinamento
comum. No caso geral, tome $s_n \nearrow f$, $t_n \nearrow g$ (\cref{thm:b3:lebesgue:approximation}): $s_n + t_n \nearrow f +
g$, e o
teorema da convergência monótona passa a aditividade ao limite. O
enunciado sobre séries é esse mesmo teorema aplicado às somas parciais.
\end{proof}

\begin{theorem}[Lema de Fatou]\label{thm:b3:lebesgue:fatou}
Para $f_n \geq 0$ \omterm{def:b3:lebesgue:measurable}{mensuráveis}:
\[
\int \liminf_n f_n \,\dd\mu \;\leq\; \liminf_n \int
f_n\,\dd\mu .
\]
\index{lema de Fatou}
\end{theorem}

\begin{proof}
Ponha $g_N = \inf_{n\geq N}f_n$: \omterm{def:b3:lebesgue:measurable}{mensurável}, $0 \leq g_N \nearrow
\liminf f_n$, e $g_N \leq f_n$ para todo $n \geq N$, de modo que
$\int g_N \leq \inf_{n \geq N}\int f_n$. Aplique a convergência monótona ao membro esquerdo: $\int\liminf f_n = \lim_N\int g_N \leq
\lim_N\inf_{n\geq N}\int f_n = \liminf\int f_n$.
\end{proof}

\begin{definition}\label{def:b3:lebesgue:l1}
Uma \omterm{def:b3:lebesgue:measurable}{função mensurável} $f \colon X \to \R$ (ou a valores em $\C$) é
\emph{integrável}\index{função integrável} se $\int\abs
f\,\dd\mu < \infty$; então $\int f = \int f^+ - \int f^-$
(partes positiva e negativa; partes real e imaginária no caso
complexo). A integral é linear nas funções integráveis (decomponha e
recombine as partes positivas; o caso complexo se reduz ao real) e
satisfaz $\abs{\int f} \leq
\int\abs f$ (caso real: $\pm\int f = \int(\pm f) \leq
\int\abs f$; caso complexo: multiplique por uma
constante unimodular para tornar a integral real). Uma propriedade vale
\emph{quase em toda parte} (q.t.p.)\index{quase em toda parte} se falha
apenas em um conjunto $\mu$-nulo; modificar $f$ em um conjunto
nulo não altera integral alguma (a diferença é dominada por $\infty\cdot\mathbf
1_N$, de
integral $0$).
\end{definition}

\begin{theorem}[Convergência dominada]\label{thm:b3:lebesgue:dct}
Seja $f_n \to f$ q.t.p., com $\abs{f_n} \leq g$ q.t.p.\ para uma $g$
\emph{\omterm{def:b3:lebesgue:l1}{integrável}} fixada. Então $f$ é \omterm{def:b3:lebesgue:l1}{integrável} e
\[
\int f_n\,\dd\mu \longrightarrow \int f\,\dd\mu,
\qquad\text{de fato}\quad \int\abs{f_n - f}\,\dd\mu \to 0 .
\]
\index{teorema da convergência dominada}
\end{theorem}

\begin{proof}
Descarte um conjunto nulo para tornar as hipóteses pontuais. $\abs f
\leq g$:
$f$ é \omterm{def:b3:lebesgue:l1}{integrável}. As funções $h_n = 2g - \abs{f_n
- f} \geq 0$ satisfazem $\liminf h_n = 2g$; Fatou dá
\[
\int 2g \leq \liminf\int\bigl(2g - \abs{f_n - f}\bigr)
= \int 2g - \limsup\int\abs{f_n - f},
\]
logo $\limsup\int\abs{f_n - f} \leq 0$ (a subtração é legítima: $\int 2g < \infty$). Por fim, $\abs{\int f_n - \int f} \leq
\int\abs{f_n - f} \to 0$.
\end{proof}

\begin{method}\label{met:b3:lebesgue:convergence}
Diante de $\lim_n\int f_n$: tente, nesta ordem --- (1) a sequência é monótona (ou
é uma série de termos não negativos)? Convergência monótona, sem
precisar de \omterm{def:b3:lebesgue:l1}{integrabilidade}. (2) Existe um único dominador \omterm{def:b3:lebesgue:l1}{integrável}
$g \geq \abs{f_n}$, obtido por estimativas grosseiras (o ``$\sup_n$'' das
estimativas)? Convergência dominada. (3) Sem dominação nem
monotonicidade? Fatou ainda limita um dos lados, e a igualdade pode
genuinamente falhar: a onda que foge, $f_n = n\mathbf
1_{\intoo0{1/n}}$, tem $\int f_n = 1$, mas $f_n \to 0$ q.t.p.
A dominação é exatamente o que impede a massa de escapar para o
infinito, na vertical ou na horizontal.
\end{method}

\section{Riemann contra Lebesgue}

\begin{theorem}[Critério de Lebesgue]
\label{thm:b3:lebesgue:riemann}
Seja $f \colon \intcc ab \to \R$ \emph{limitada}. Então $f$ é \omterm{def:b3:lebesgue:l1}{Riemann-integrável} se, e
somente se, $f$ é \omterm{def:b3:topology:continuity}{contínua} $\lambda$-quase em toda parte; nesse caso,
$f$ é \omterm{def:b3:lebesgue:l1}{Lebesgue-integrável} e as duas integrais coincidem.\index{integral
de Riemann contra a de Lebesgue}
\end{theorem}

\begin{proof}
Para uma subdivisão $\sigma = (a = x_0 < \dots < x_N = b)$, sejam $L_\sigma$ e $U_\sigma$ as funções escada iguais,
em cada $\intoo{x_{i-1}}{x_i}$, a $m_i = \inf_{[x_{i-1}, x_i]}f$ e a $M_i = \sup$; as somas de Darboux são suas integrais
(Riemann e Lebesgue coincidem em funções escada, ambas dando $\sum
m_i\Delta x_i$).
Tome uma sequência de subdivisões $\sigma_n$, cada uma refinando a anterior,
de malha $\to 0$, com somas de Darboux convergindo às integrais de
Darboux inferior e superior de $f$. Os refinamentos tornam $L_{\sigma_n}$ não
decrescente e $U_{\sigma_n}$ não crescente pontualmente fora do conjunto
enumerável $D$ de todos os pontos de divisão; chame os limites de
$\ell$ e $u$ (\omterm{def:b3:lebesgue:measurable}{mensuráveis}, \cref{prop:b3:lebesgue:stability}). Para $x \notin
D$,
escrevendo $I_n(x)$ para o intervalo aberto de $\sigma_n$ que contém
$x$: $\ell(x) = \sup_n\inf_{I_n(x)}f$ e $u(x) =
\inf_n\sup_{I_n(x)}f$; como as malhas encolhem a $0$, essas são as
\emph{envoltórias} inferior e superior de $f$ em $x$ --- $u(x) - \ell(x)$ é
a oscilação de $f$ em $x$ ---, de modo que \emph{$\ell(x) = u(x)$ se, e
somente se, $f$ é \omterm{def:b3:topology:continuity}{contínua} em $x$}. Por convergência monótona ou
dominada (função limitada, intervalo finito):
\[
\int_{\intcc ab}\ell\,\dd\lambda = \lim_n\int L_{\sigma_n}
= \underline{\int}f,
\qquad
\int_{\intcc ab}u\,\dd\lambda = \overline{\int}f .
\]
$f$ \omterm{def:b3:lebesgue:l1}{Riemann-integrável} $\iff$ $\underline\int f =
\overline\int f$ $\iff$ $\int(u - \ell) = 0$ $\iff$
$u = \ell$ q.t.p.\ ($u - \ell \geq 0$; \cref{exo:b3:lebesgue:5}) $\iff$ $f$ \omterm{def:b3:topology:continuity}{contínua} q.t.p.
Nesse caso, $\ell \leq f \leq u$ com $\ell =
u$ q.t.p.: $f$ é igual q.t.p.\ à \omterm{def:b3:lebesgue:measurable}{função
mensurável} $\ell$ e, portanto, é \omterm{def:b3:lebesgue:measurable}{Lebesgue-mensurável} (\omterm{def:b3:complete:complete}{completude} de
$\lambda$), com $\int f\,\dd\lambda = \int\ell\,\dd\lambda = \underline\int f =
\int_a^bf$.
\end{proof}

\begin{example}\label{ex:b3:lebesgue:riemannexamples}
$\mathbf 1_\Q$ não é \omterm{def:b3:topology:continuity}{contínua} em ponto algum: não é \omterm{def:b3:lebesgue:l1}{Riemann-integrável} --- mas é
trivial para Lebesgue: $\int\mathbf 1_\Q\,\dd\lambda =
\lambda(\Q) = 0$. A função de Thomae (\,$\frac1q$ nos racionais
$\frac pq$ e $0$ nos demais pontos) é \omterm{def:b3:topology:continuity}{contínua} exatamente nos
irracionais: é \omterm{def:b3:lebesgue:l1}{Riemann-integrável} com integral $0$. E as integrais
de Riemann \emph{impróprias} são uma noção diferente: $\int_0^\infty\frac{\sin x}x\,\dd x$ converge
como limite de $\int_0^A$ (o problema de fim de semana a calcula: $= \frac\pi2$), mas
$\frac{\sin x}x \notin L^1(\intoo0{+\infty})$: a integral absoluta diverge como a série harmônica
(\cref{exo:b3:lebesgue:6}). A teoria de Lebesgue troca a convergência
condicional por teoremas de limite robustos.
\end{example}

\section{Integrais com parâmetros}

Ao longo de toda esta seção, $(X, \mathcal A, \mu)$ é um \omterm{def:b3:measure:measure}{espaço de medida}, $T$ é um
espaço métrico (o parâmetro) e $f \colon T \times X \to \C$ é tal que $f(t, \cdot)$ é \omterm{def:b3:lebesgue:l1}{integrável} para
cada $t$; ponha $F(t) = \int_X
f(t, x)\,\dd\mu(x)$.

\begin{theorem}[Continuidade]\label{thm:b3:lebesgue:paramcont}
Suponha: $t \mapsto f(t,x)$ é \omterm{def:b3:topology:continuity}{contínua} em $t_0$ para q.t.p.\ $x$, e existe
uma \omterm{def:b3:lebesgue:l1}{função integrável} $g$ com $\abs{f(t,x)} \leq
g(x)$ para todo $t$ em uma
\omterm{def:b3:topology:topology}{vizinhança} de $t_0$ e q.t.p.\ $x$. Então $F$ é \omterm{def:b3:topology:continuity}{contínua} em
$t_0$.\index{integrais com parâmetro}
\end{theorem}

\begin{proof}
Para uma sequência qualquer $t_n \to t_0$: $f(t_n, \cdot) \to f(t_0,
\cdot)$ q.t.p., dominada por $g$:
a convergência dominada dá $F(t_n) \to F(t_0)$; e a \omterm{def:b3:topology:continuity}{continuidade} sequencial basta em
espaços métricos (\cref{rem:b3:topology:sequences}).
\end{proof}

\begin{theorem}[Derivação sob o sinal de integral]
\label{thm:b3:lebesgue:paramdiff}
Seja $T$ um intervalo aberto de $\R$. Suponha: para q.t.p.\
$x$, $t \mapsto f(t,x)$ é derivável em $T$, com
\[
\Bigl|\frac{\partial f}{\partial t}(t, x)\Bigr| \leq g(x)
\quad \text{para todo } t \in T,\ \text{q.t.p. } x,
\]
e $g$ \omterm{def:b3:lebesgue:l1}{integrável}. Então $F$ é derivável em $T$ com $F'(t)
= \int_X \frac{\partial f}{\partial t}(t, x)\,\dd\mu(x)$.
\end{theorem}

\begin{proof}
Fixe $t$ e $h_n \to 0$: os quocientes de diferenças
\[
\varphi_n(x) = \frac{f(t + h_n, x) - f(t, x)}{h_n}
\longrightarrow \frac{\partial f}{\partial t}(t,x)
\quad\text{q.t.p.},
\]
e a desigualdade do valor médio limita $\abs{\varphi_n(x)} \leq
\sup_{s}\abs{\partial_tf(s,x)} \leq g(x)$: a convergência dominada
se aplica, e $\frac{F(t + h_n) - F(t)}{h_n} = \int\varphi_n \to
\int\partial_t f(t, \cdot)$.
\end{proof}

\begin{example}[A função Gama]\label{ex:b3:lebesgue:gamma}
Para $t > 0$, ponha\index{função Gama}
\[
\Gamma(t) = \int_0^{+\infty} x^{t-1}\eu^{-x}\,\dd x .
\]
A integral converge: perto de $0$, $x^{t-1}$ é \omterm{def:b3:lebesgue:l1}{integrável} ($t >
0$);
no infinito, $x^{t-1}\eu^{-x} \leq C\eu^{-x/2}$. A integração por partes (em $[\varepsilon, A]$ e, depois,
passando ao limite por convergência monótona) dá a equação funcional
$\Gamma(t + 1) =
t\,\Gamma(t)$, donde $\Gamma(n+1) = n!$: o fatorial interpolado. Em todo $\intcc ab \subseteq \intoo0{+\infty}$, $\partial_t\bigl(x^{t-1}\eu^{-x}\bigr) = \ln
x\cdot x^{t-1}\eu^{-x}$ é
dominada por $\abs{\ln x}(x^{a-1} +
x^{b-1})\eu^{-x}$, \omterm{def:b3:lebesgue:l1}{integrável}: $\Gamma$ é $\mathcal C^1$ e, por indução,
$\mathcal C^\infty$, com $\Gamma^{(k)}(t) =
\int_0^\infty(\ln x)^kx^{t-1}\eu^{-x}\dd x$. O valor $\Gamma(\frac12) = \sqrt\pi$ é a integral gaussiana disfarçada
(\cref{pb:b3:lebesgue:1}).
\end{example}

\begin{omfigure}
\begin{tikzpicture}
\begin{axis}[omaxis, width=10.6cm, height=5.4cm,
    xmin=0, xmax=25, ymin=-0.25, ymax=1.05,
    xtick={0, 3.1416, 6.2832, 9.4248, 12.566, 15.708, 18.850,
      21.991, 25.133},
    xticklabels={$0$, $\pi$, $2\pi$, $3\pi$, $4\pi$, $5\pi$,
      $6\pi$, $7\pi$, $8\pi$},
    ytick={0, 1}]
  \addplot[omDef!40] coordinates {(0,0) (25,0)};
  \addplot[omThm, thick, domain=0.02:25, samples=600]
    {sin(deg(x))/x};
  \addplot[omDef, dashed, domain=0.02:25, samples=200]
    {1/x};
  \addplot[omDef, dashed, domain=0.02:25, samples=200]
    {-1/x};
\end{axis}
\end{tikzpicture}

{\small O integrando $\frac{\sin x}x$: a integral imprópria $\int_0^\infty$ converge por
cancelamento alternado entre os arcos, mas as áreas $\abs{\cdot}$ dos arcos se
comportam como $\frac2{\pi k}$ --- uma série harmônica: $\frac{\sin x}x
\notin L^1$. Ser
\omterm{def:b3:lebesgue:l1}{Lebesgue-integrável} é ser absolutamente \omterm{def:b3:lebesgue:l1}{integrável}.}
\end{omfigure}

\section{Exercícios}

\begin{exercise}[$\star$]\label{exo:b3:lebesgue:1}
(a) Mostre que uma função monótona $\R \to \R$ é \omterm{def:b3:lebesgue:measurable}{Borel-mensurável}, e que uma
derivada (de uma função derivável em toda parte) é \omterm{def:b3:lebesgue:measurable}{Borel-mensurável}.
(b) Mostre que $f \colon X \to \R$ é \omterm{def:b3:lebesgue:measurable}{mensurável} se, e somente se, $\{f > q\}
\in \mathcal A$ para todo
$q$ \emph{racional}.
\end{exercise}

\begin{exercise}[$\star$]\label{exo:b3:lebesgue:2}
Calcule, com justificativa \omterm{def:b3:complete:complete}{completa}:
\[
\lim_{n\to\infty}\int_0^{+\infty}
\frac{\cos x}{(1 + x/n)^{n}}\,\dd x,
\qquad
\lim_{n\to\infty}\int_0^1 \frac{n\,x^{n-1}}{1 + x}\,\dd x .
\]
\emph{(Para o segundo: substitua $u = x^n$ antes de dominar.)}
\end{exercise}

\begin{exercise}[$\star\star$]\label{exo:b3:lebesgue:3}
(a) Exiba desigualdade estrita no lema de Fatou.
(b) Exiba $f_n \to 0$ pontualmente com $\int f_n = 1$ de três maneiras: fuga em
altura, em largura e para o infinito. Qual hipótese isolada da
convergência dominada cada uma delas viola?
(c) Mostre que, no lema de Fatou, não se pode trocar $\liminf$ por $\limsup$ em
nenhum dos dois lados.
\end{exercise}

\begin{exercise}[$\star\star$]\label{exo:b3:lebesgue:4}
(a) Mostre que $\displaystyle\int_0^{+\infty}\frac{x}{\eu^x -
1}\,\dd x = \sum_{n\geq1}\frac1{n^2} = \frac{\pi^2}6$ \emph{(expanda $\frac1{\eu^x - 1}$ em série geométrica e integre
termo a termo --- que teorema o permite?)}.
(b) (O sonho do segundanista) Mostre que $\displaystyle\int_0^1 x^{-x}\,\dd x = \sum_{n\geq1}n^{-n}$. \emph{(Escreva $x^{-x} = \eu^{-x\ln x} = \sum_k\frac{(-x\ln
x)^k}{k!}$ e
calcule $\int_0^1(-x\ln x)^k\dd x$ substituindo $x = \eu^{-u/(k+1)}$, reconhecendo $\Gamma$.)}
\end{exercise}

\begin{exercise}[$\star\star$]\label{exo:b3:lebesgue:5}
(a) Mostre que $f \geq 0$ \omterm{def:b3:lebesgue:measurable}{mensurável} com $\int f\,\dd\mu = 0$ satisfaz $f = 0$ q.t.p.
\emph{(Considere $\{f \geq 1/n\}$ e a desigualdade de Markov: $\mu(\{f \geq a\}) \leq \frac1a\int f$ ---
demonstre-a.)}
(b) Mostre que uma $f$ \omterm{def:b3:lebesgue:l1}{integrável} é finita q.t.p.
(c) Mostre que, se $\int_A f\,\dd\mu = 0$ para \emph{todo} $A$ \omterm{def:b3:lebesgue:measurable}{mensurável}, então
$f = 0$ q.t.p.
\end{exercise}

\begin{exercise}[$\star\star$]\label{exo:b3:lebesgue:6}
(a) Aplique o~\cref{thm:b3:lebesgue:riemann} para decidir a \omterm{def:b3:lebesgue:l1}{integrabilidade} de Riemann
de: $\mathbf 1_\Q$; a função de Thomae; $\mathbf
1_K$ para $K$ um conjunto de Cantor
gordo (\cref{exo:b3:measure:5}).
(b) Mostre que $\int_1^{+\infty}\abs{\frac{\sin x}x}\,\dd x =
+\infty$, ao passo que $\lim_{A\to\infty}\int_1^A\frac{\sin
x}x\,\dd x$ existe \emph{(integre por
partes)}: convergência imprópria sem \omterm{def:b3:lebesgue:l1}{integrabilidade}.
\end{exercise}

\begin{exercise}[$\star\star$]\label{exo:b3:lebesgue:7}
Justifique que $F(t) = \int_0^{+\infty}\eu^{-x^2}\cos(tx)\,\dd x$ é $\mathcal C^1$ em $\R$ e satisfaz $F'(t) = -\frac t2F(t)$ \emph{(integre
por partes)}; deduza $F(t) =
F(0)\,\eu^{-t^2/4}$. (Com o $F(0) = \frac{\sqrt\pi}2$ obtido no problema de fim de
semana: a gaussiana é essencialmente sua própria transformada de
Fourier --- o~\cref{ch:b3:fouriertransform} sistematizará isso.)
\end{exercise}

\begin{exercise}[$\star\star\star$]\label{exo:b3:lebesgue:8}
(Frullani) Sejam $0 < a < b$. Mostre que
\[
\int_0^{+\infty}\frac{\eu^{-ax} - \eu^{-bx}}{x}\,\dd x =
\ln\frac ba,
\]
escrevendo o integrando como $\int_a^b \eu^{-xt}\,\dd t$ e justificando a troca pela teoria
não negativa (\cref{cor:b3:lebesgue:additivity} em forma \omterm{def:b3:topology:continuity}{contínua} --- antecipe Tonelli, ou
fatie $[a,b]$ em $n$ partes iguais e passe ao limite).
\end{exercise}

\begin{exercise}[$\star\star$]\label{exo:b3:lebesgue:9}
Seja $f \geq 0$ \omterm{def:b3:lebesgue:measurable}{mensurável} em $(X, \mathcal A, \mu)$. Mostre que $\nu(A) = \int_A f\,\dd\mu$ define uma \omterm{def:b3:measure:measure}{medida} (de
\emph{densidade}\index{densidade de uma medida} $f$ em relação a
$\mu$) e que $\int g\,\dd\nu = \int gf\,\dd\mu$ para toda $g \geq 0$ \omterm{def:b3:lebesgue:measurable}{mensurável} \emph{(demonstre-o
para indicadoras, depois para \omterm{def:b3:lebesgue:simple}{funções simples} e, então, por
convergência monótona --- a máquina padrão)}.
\end{exercise}

\begin{exercise}[$\star\star\star$]\label{exo:b3:lebesgue:10}
(Um fracasso ao estilo de Weierstrass) Defina $f(t) =
\int_0^{+\infty}\frac{\sin(tx)}{x(1 + x^2)}\,\dd x$.
(a) Mostre que $f$ está bem definida e é \omterm{def:b3:topology:continuity}{contínua} em $\R$, e que
$\mathcal C^1$ com $f'(t) = \int_0^\infty\frac{\cos(tx)}{1 +
x^2}\dd x$ para todo $t$ --- mas que derivar \emph{de novo} sob
o sinal de integral é ilegítimo.
(b) Admitindo $f'(t) = \frac\pi2\eu^{-t}$ para $t > 0$ (demonstrado no~\cref{ch:b3:residues}),
quanto vale $\int_0^\infty\frac{x\sin(tx)}{1+x^2}\dd x$ para $t > 0$, e por que sua fórmula confirma o
fracasso de (a)?
\end{exercise}

\begin{exercise}[$\star\star$]\label{exo:b3:lebesgue:11}
(Lema de Scheffé) Sejam $f_n, f \geq 0$ \omterm{def:b3:lebesgue:l1}{integráveis} com $f_n \to f$ q.t.p.\ e $\int f_n \to \int f$.
(a) Mostre que $\int\abs{f_n - f} \to 0$. \emph{(Aplique a convergência dominada a $g_n = (f - f_n)^+ \leq f$ e
escreva $\int\abs{f_n - f} = 2\int g_n - \int(f - f_n)$.)}
(b) Mostre com um exemplo que a hipótese $\int f_n \to \int
f$ não pode ser descartada
(uma onda deslizante ou que se concentra), e que a conclusão falha para
$f_n$ com sinal, sem controle em valor absoluto: $f_n = n\mathbf 1_{\intoc0{1/n}} -
n\mathbf 1_{\intoc{-1/n}0}$ tem $f_n \to 0$
q.t.p., $\int f_n
= 0 \to 0$, e no entanto $\int\abs{f_n} = 2$.
(c) Aplicação (\omterm{ex:b3:lebesgue:gamma}{densidades}): se \omterm{ex:b3:lebesgue:gamma}{densidades} de probabilidade $p_n
\to p$
q.t.p., então automaticamente $\int\abs{p_n - p} \to 0$: a convergência pontual de
\omterm{ex:b3:lebesgue:gamma}{densidades} é convergência $L^1$ --- um ganho de convergência de
graça.
\end{exercise}

\begin{exercise}[$\star\star$]\label{exo:b3:lebesgue:12}
Limites clássicos, com justificativa \omterm{def:b3:complete:complete}{completa} por convergência monótona
ou dominada:
\[
\text{(a)}\ \lim_{n\to\infty}\int_0^n\Bigl(1 -
\frac xn\Bigr)^n\eu^{x/2}\,\dd x, \qquad
\text{(b)}\ \lim_{n\to\infty}\int_0^1\frac{n\,x^{n-1}}{1 +
x}\,\dd x,
\]
\[
\text{(c)}\ \lim_{n\to\infty}\int_0^\infty
\frac{\dd x}{(1 + x/n)^n\,x^{1/n}} .
\]
\emph{(Para (a): $(1 - x/n)^n \nearrow \eu^{-x}$ com $x$ fixado --- demonstre a monotonicidade
via $\log$; para (b), integre por partes ou substitua $x = u^{1/n}$ e
identifique uma concentração na fronteira; para (c), encontre um
dominador \omterm{def:b3:lebesgue:l1}{integrável} válido para todo $n \geq 2$ separando em $x =
1$.)}
\end{exercise}

\section{Problema: duas integrais célebres}

\begin{problem}[{Problema de fim de semana --- a integral gaussiana e a
integral de Dirichlet, só com parâmetros}]
\label{pb:b3:lebesgue:1}
Duas integrais governam a análise aplicada:
\[
G = \int_{-\infty}^{+\infty}\eu^{-x^2}\dd x = \sqrt\pi,
\qquad
D = \int_0^{+\infty}\frac{\sin x}{x}\,\dd x = \frac\pi2
\]
(a segunda como integral imprópria, \cref{ex:b3:lebesgue:riemannexamples}). Demonstramos ambas
usando apenas as ferramentas deste capítulo.

\medskip
\noindent\textbf{Parte I --- A gaussiana.} Para $t \geq 0$, ponha
\[
A(t) = \Bigl(\int_0^t\eu^{-x^2}\dd x\Bigr)^{2},
\qquad
B(t) = \int_0^1\frac{\eu^{-t^2(1 + x^2)}}{1 + x^2}\,\dd x .
\]
\begin{enumerate}
  \item Justifique que $A$ e $B$ são $\mathcal C^1$ em $\intoo0{+\infty}$ e calcule
        $A'$ e $B'$; mostre que $A'(t) + B'(t) = 0$. \emph{(Em $B'$,
        substitua $u =
        tx$.)}
  \item Calcule $A(0) + B(0)$ e $\lim_{t\to+\infty}(A + B)(t)$ --- justifique a passagem ao limite sob a
        integral em $B$.
  \item Conclua $\int_0^\infty \eu^{-x^2}\dd x =
        \frac{\sqrt\pi}2$, logo $G = \sqrt\pi$, e deduza $\Gamma(\tfrac12) = \sqrt\pi$ \emph{(substitua $x =
        u^2$
        em $\Gamma(\frac12)$)}.
\end{enumerate}

\medskip
\noindent\textbf{Parte II --- A integral de Dirichlet.} Para $t
\geq 0$,
ponha
\[
F(t) = \int_0^{+\infty}\eu^{-tx}\,\frac{\sin x}{x}\,\dd x .
\]
\begin{enumerate}[resume]
  \item Mostre que a integral que define $F(t)$ converge para todo
        $t > 0$ como integral de Lebesgue e, para $t = 0$, como
        integral imprópria; mostre que $D =
        \lim_{A\to\infty}\int_0^A\frac{\sin x}x\dd x$ existe \emph{(integre por
        partes em $[\pi, A]$)}.
  \item Mostre que $F$ é $\mathcal C^1$ em $\intoo0{+\infty}$ com
        \[
        F'(t) = -\int_0^{+\infty}\eu^{-tx}\sin x\,\dd x
        = -\frac{1}{1 + t^2}
        \]
        \emph{(dominação em $[t_0, \infty)$ para cada $t_0 >
        0$; a última integral por
        duas integrações por partes ou por exponenciais complexas)}.
  \item Mostre que $F(t) \to 0$ quando $t \to +\infty$, e deduza $F(t)
        = \frac\pi2 - \arctan t$ em $\intoo0{+\infty}$.
  \item O ponto delicado: $D = \lim_{t\to0^+}F(t)$. Demonstre-o por controle uniforme da
        cauda: para $0 \leq t \leq 1$ e $A \geq \pi$, integre por partes para mostrar que
        \[
        \Bigl|\int_A^{+\infty}\eu^{-tx}\frac{\sin
        x}x\,\dd x\Bigr| \leq \frac{C}{A}
        \]
        com $C$ independente de $t$ \emph{(derive $\frac{\eu^{-tx}}x$ e limite
        $\abs{\cos}$ por $1$; note que $t\eu^{-tx} \leq 1/x\cdot(tx\eu^{-tx})$ com $\sup_{u\geq0}u\eu^{-u} < 1$)}; depois, separe
        $F(t) - D$ em $[0, A]$ (onde a convergência dominada se aplica
        quando $t \to 0$) e $[A, \infty)$.
  \item Conclua: $D = \frac\pi2$.
\end{enumerate}

\medskip
\noindent\textbf{Parte III --- Dividendos.}
\begin{enumerate}[resume]
  \item Calcule $\int_0^{+\infty}\frac{\sin^2x}{x^2}\,\dd x$ \emph{(integre por partes e reduza a $D$ via
        $\sin
        2x = 2\sin x\cos x$)}.
  \item Calcule $\int_0^{+\infty}\frac{1 -
        \cos x}{x^2}\,\dd x$ e confira a coerência dos dois resultados.
  \item Para $a > 0$, calcule $\int_0^{+\infty}\frac{\sin(ax)}x\dd x$ e $\int_{-\infty}^{+\infty}\eu^{-ax^2}\dd x$, e registre as regras de
        reescalamento (elas serão os cavalos de batalha do~\cref{ch:b3:fouriertransform}).
  \item Explique com precisão por que $D$ não poderia ter sido
        tratada diretamente por convergência dominada em $t = 0$
        (não há dominador \omterm{def:b3:lebesgue:l1}{integrável} em $[0,1]\times[0,\infty)$), e por que a separação de
        cauda da questão 7 é o substituto honesto --- esse padrão
        (``\omterm{def:b3:lebesgue:l1}{integrabilidade} uniforme das caudas'') reaparece por toda a
        análise.
\end{enumerate}

\medskip
\noindent\textbf{Parte IV --- A \omterm{ex:b3:lebesgue:gamma}{função Gama} segundo Bohr e Mollerup.} A
função $\Gamma$ (\cref{ex:b3:lebesgue:gamma}) satisfaz $\Gamma(1) = 1$ e $\Gamma(x+1) = x\Gamma(x)$ --- mas
infinitas outras funções também satisfazem (multiplique por qualquer
ondulação $1$-periódica). Uma única condição de convexidade fixa
$\Gamma$ de maneira única, e suas identidades mais profundas caem
então mecanicamente. Uma função positiva $f$ em um intervalo é
\emph{log-convexa} se $\log f$ é convexa.
\begin{enumerate}[resume]
  \item Mostre que log-convexa implica convexa, que produtos de funções
        log-convexas e suas composições com aplicações afins são
        log-convexos e --- via a desigualdade de Hölder para duas
        funções $\int\abs{uv} \leq
        \bigl(\int\abs u^p\bigr)^{1/p}\bigl(\int\abs
        v^q\bigr)^{1/q}$, demonstrada diretamente a partir da desigualdade
        de Young --- que $\Gamma$ é log-convexa em $\intoo0\infty$.
  \item (Lema das inclinações) Seja $g$ convexa em $\intoo0\infty$ com $g(n+1) - g(n) = \log n$
        para todo inteiro $n
        \geq 1$. Para $x \in \intoc01$ e $n \geq 2$, compare as
        inclinações de $g$ em $[n-1, n]$, $[n, n+x]$ e $[n, n+1]$, e deduza
        \[
        x\log(n-1) \;\leq\; g(n + x) - g(n) \;\leq\; x\log n .
        \]
  \item (Bohr--Mollerup) Seja $f > 0$ com $f(1) = 1$, $f(x+1) = xf(x)$ e $\log f$
        convexa. Desenrolando a recursão em $f(n + x) = x(x+1)\cdots(x + n -
        1)\,f(x)$ e $f(n) = (n-1)!$, deduza da
        questão 14 que, para $x \in \intoc01$,
        \[
        f(x) = \lim_{n\to\infty}
        \frac{n!\,n^x}{x(x+1)\cdots(x+n)} :
        \]
        $f$ é única, logo $f = \Gamma$, e vale a \emph{fórmula-limite de
        Gauss} (estenda a todo $x > 0$ pela recursão).
  \item Defina a \emph{função Beta} $B(x, y) =
        \int_0^1t^{x-1}(1-t)^{y-1}\,\dd t$ ($x, y > 0$). Demonstre a
        convergência, a recursão $B(x+1, y) =
        \frac{x}{x+y}\,B(x, y)$ (integre por partes) e $B(1, y) = \frac1y$.
  \item Mostre que $x \mapsto B(x, y)$ é log-convexa (Hölder de novo) e aplique
        Bohr--Mollerup a
        \[
        f(x) = \frac{B(x, y)\,\Gamma(x + y)}{\Gamma(y)}
        \]
        para concluir a \emph{fórmula de Euler}: $B(x, y) =
        \dfrac{\Gamma(x)\Gamma(y)}{\Gamma(x+y)}$ --- sem nenhuma
        integral dupla.
  \item Calcule $B(\frac12, \frac12)$ diretamente (substitua $t
        = \sin^2\theta$) e deduza $\Gamma(\frac12) =
        \sqrt\pi$: a
        integral gaussiana da Parte I, recuperada por convexidade pura.
        Compare as duas demonstrações em uma frase cada.
  \item (Duplicação de Legendre) Mostre que
        \[
        g(x) = \frac{2^{x-1}}{\sqrt\pi}\,
        \Gamma\Bigl(\frac x2\Bigr)
        \Gamma\Bigl(\frac{x+1}2\Bigr)
        \]
        satisfaz as três hipóteses de Bohr--Mollerup e conclua $g = \Gamma$,
        isto é, $\Gamma(2z) =
        \frac{2^{2z-1}}{\sqrt\pi}\,\Gamma(z)\,\Gamma(z +
        \tfrac12)$ para todo $z > 0$.
  \item Deduza a forma fechada $\Gamma\bigl(n + \tfrac12\bigr)
        = \dfrac{(2n)!}{4^n\,n!}\sqrt\pi$ e demonstre, pelo lema das
        inclinações aplicado a $\log\Gamma$ em torno de inteiros grandes, as
        assintóticas
        \[
        \frac{\Gamma(n + \frac12)}{\Gamma(n)\,\sqrt n}
        \longrightarrow 1 .
        \]
  \item Combine as duas últimas questões na assintótica do coeficiente
        binomial central
        \[
        \binom{2n}{n} \sim \frac{4^n}{\sqrt{\pi n}},
        \]
        e verifique numericamente para $n = 10$ ($\binom{20}{10} =
        184756$, contra
        $4^{10}/\sqrt{10\pi} \approx
        187079$: razão $\approx 0.988$).
  \item (Síntese) A constante $\sqrt\pi$ já apareceu como a integral
        gaussiana (Parte I), como $B(\frac12,
        \frac12)$ (questão 18) e dentro da
        duplicação (questão 19); a estimativa do binomial central
        antecipa tanto Stirling (o problema de fim de semana do~\cref{ch:b3:product}) quanto de Moivre--Laplace. Mapeie as conexões:
        que enunciados são equivalentes a quais, e o que cada técnica
        --- derivação sob a integral contra convexidade --- traz que a
        outra não pode trazer?
\end{enumerate}

\medskip
\noindent\textbf{Parte V --- Mais três dividendos.}
\begin{enumerate}[resume]
  \item (Wallis, via Beta) Para $p > -1$, substitua $t =
        \sin^2\theta$ para
        mostrar
        \[
        W_p = \int_0^{\pi/2}\sin^p\theta\,\dd\theta
        = \frac12\,B\Bigl(\frac{p+1}2, \frac12\Bigr),
        \]
        e deduza da recursão da Beta (questão 16) que $W_{n+2} = \frac{n+1}{n+2}\,W_n$ para
        inteiros $n \geq
        0$. Calcule $W_{2n}$ e $W_{2n+1}$ em forma fechada,
        mostre $W_{2n+1}/W_{2n} \to 1$ por confinamento e conclua com o \emph{produto de
        Wallis}
        \[
        \frac\pi2 = \lim_{n\to\infty}
        \prod_{k=1}^{n}\frac{4k^2}{4k^2 - 1} .
        \]
  \item (A gaussiana encontra uma frequência) Para $b \in \R$, ponha
        \[
        \Phi(b) = \int_{-\infty}^{+\infty}
        \eu^{-x^2}\cos(2bx)\,\dd x .
        \]
        Mostre que $\Phi$ é $\mathcal C^1$ em $\R$, que uma integração por
        partes fornece a equação diferencial $\Phi'(b) = -2b\,\Phi(b)$ e conclua
        \[
        \Phi(b) = \sqrt\pi\,\eu^{-b^2} :
        \]
        a gaussiana se reproduz sob essa transformada --- a única
        identidade sobre a qual o~\cref{ch:b3:fouriertransform} vai funcionar.
  \item (A integral de Frullani) Para $0 < a < b$, mostre que
        \[
        \int_0^{+\infty}
        \frac{\eu^{-ax} - \eu^{-bx}}{x}\,\dd x
        = \log\frac ba ,
        \]
        derivando no parâmetro $a$ (justifique a dominação em todo
        $\intco{a_0}{+\infty}$, $a_0 > 0$, e identifique a constante fazendo $a \to b$). Onde
        exatamente o integrando precisa de sua singularidade removível
        em $x = 0$?
\end{enumerate}
\end{problem}
